\documentclass[11pt,a4paper]{article}

% =========================
% PACKAGES
% =========================
\usepackage[utf8]{inputenc}
\usepackage[T1]{fontenc}
\usepackage[french]{babel}
\usepackage{lmodern}
\usepackage{geometry}
\geometry{margin=1.6cm}
\usepackage{xcolor}
\usepackage{amsmath,amssymb}
\usepackage{tikz}
\usepackage{array}
\usepackage{enumitem}
\usepackage{fancyhdr}
\usepackage{multicol}
\usepackage{tcolorbox}

% =========================
% STYLE
% =========================
\pagestyle{fancy}
\fancyhf{}
\lhead{\textbf{Mathématiques -- Sixième}}
\rhead{\textbf{DS -- Symétrie axiale}}
\cfoot{\thepage}

\setlength{\parindent}{0pt}
\setlength{\parskip}{0.35em}

\definecolor{bleu}{RGB}{30,80,160}
\definecolor{grisclair}{RGB}{245,245,245}

\newcommand{\points}[1]{\hfill \textbf{(#1 points)}}

\newtcolorbox{consignes}{
  colback=grisclair,
  colframe=bleu,
  title=\textbf{Consignes},
  fonttitle=\bfseries
}

\begin{document}

% =========================
% TITRE
% =========================
\begin{center}
    {\Large \textbf{Devoir surveillé de mathématiques}}\\[0.2cm]
    {\large \textbf{Chapitre : Symétrie axiale}}\\[0.2cm]
    Niveau sixième
\end{center}

\vspace{0.2cm}

\begin{tabular}{|p{0.45\linewidth}|p{0.45\linewidth}|}
\hline
Nom : & Prénom : \\
\hline
Classe : & Date : \\
\hline
\end{tabular}

\vspace{0.4cm}

% =========================
% TITRE
% =========================
\begin{center}
    {\Large \textbf{Devoir maison de mathématiques}}\\[0.2cm]
    {\large \textbf{Chapitre : Symétrie axiale}}\\[0.2cm]
    Niveau sixième
\end{center}

\vspace{0.2cm}

\begin{tabular}{|p{0.45\linewidth}|p{0.45\linewidth}|}
\hline
Nom : & Prénom : \\
\hline
Classe : & Date : \\
\hline
\end{tabular}

\vspace{0.4cm}

\begin{consignes}
\begin{itemize}[label=\textbullet]
    \item Le devoir est à faire proprement, au crayon à papier pour les constructions.
    \item Les traits de construction doivent rester visibles.
    \item On utilisera la règle graduée, l'équerre et le compas si nécessaire.
    \item Toute réponse doit être justifiée lorsque cela est demandé.
    \item pour ce DM vous pouvez répondre sur le sujet comme j'ai laissé la place.
    \item les exercices à faire en priorités sont les exercices 1,3, 4
\end{itemize}
\end{consignes}

\vspace{0.3cm}

\begin{center}
\textbf{Barème : 31 points}
\end{center}

% ============================================================
\section*{Exercice 1 -- Vocabulaire et reconnaissance \points{4}}

On rappelle qu'une figure possède un axe de symétrie si, lorsqu'on la plie suivant cet axe, les deux parties se superposent exactement.

\medskip

\textbf{1.} Pour chaque figure ci-dessous, dire si la droite en pointillés est un axe de symétrie.  
Écrire \textbf{oui} ou \textbf{non} sous chaque figure. \points{2}

\begin{center}
\begin{tikzpicture}[scale=0.9]

% Figure A : rectangle avec axe vertical
\draw[thick] (0,0) rectangle (2,1.4);
\draw[dashed, thick] (1,-0.3)--(1,1.7);
\node at (1,-0.7) {Figure A : \dotfill};

% Figure B : triangle quelconque avec faux axe
\draw[thick] (4,0)--(6,0)--(5.4,1.5)--cycle;
\draw[dashed, thick] (5,-0.3)--(5,1.8);
\node at (5,-0.7) {Figure B : \dotfill};

% Figure C : cercle avec axe
\draw[thick] (9,0.7) circle (0.8);
\draw[dashed, thick] (8.2,0.7)--(9.8,0.7);
\node at (9,-0.7) {Figure C : \dotfill};

\end{tikzpicture}
\end{center}

\vspace{0.4cm}

\textbf{2.} Compléter les phrases suivantes avec les mots :  
\textbf{perpendiculaire}, \textbf{même distance}, \textbf{milieu}. \points{2}

\begin{enumerate}[label=\alph*)]
    \item Si deux points $A$ et $A'$ sont symétriques par rapport à une droite $(d)$, alors $A$ et $A'$ sont à la ........................................ de la droite $(d)$.
    \item La droite $(d)$ est \dotfill au segment $[AA']$.
    \item La droite $(d)$ passe par le \dotfill du segment $[AA']$.
\end{enumerate}

% ============================================================
\section*{Exercice 2 -- Construire le symétrique de points \points{5}}

Dans la figure ci-dessous, construire les points $A'$, $B'$, $C'$ et $D'$, symétriques respectifs des points $A$, $B$, $C$ et $D$ par rapport à la droite $(d)$.

\medskip

\textbf{Attention :} on laissera visibles les traits de construction.

\begin{center}
\begin{tikzpicture}[scale=0.75]

% grille
\draw[step=1cm, gray!35, very thin] (-5,-3) grid (5,5);

% axe
\draw[very thick, bleu] (0,-3)--(0,5);
\node[bleu, above] at (0,5) {$(d)$};

% points
\filldraw (-3,4) circle (2pt) node[above left] {$A$};
\filldraw (-2,1) circle (2pt) node[above left] {$B$};
\filldraw (-4,-1) circle (2pt) node[below left] {$C$};
\filldraw (-1,-2) circle (2pt) node[below left] {$D$};

% graduations invisibles? non
\end{tikzpicture}
\end{center}

\textbf{Questions :}

\begin{enumerate}[label=\arabic*)]
    \item Construire $A'$, $B'$, $C'$ et $D'$. \points{4}
    \item Que peut-on dire de la droite $(d)$ par rapport aux segments $[AA']$, $[BB']$, $[CC']$ et $[DD']$ ? \points{1}
\end{enumerate}

\vspace{1cm}

Réponse à la question 2 :

\dotfill

\dotfill

% ============================================================
\section*{Exercice 3 -- Construire le symétrique d'une figure \points{6}}

On a tracé une figure composée de plusieurs segments.  
Construire son symétrique par rapport à la droite $(d)$.

\medskip

\textbf{Méthode conseillée :}
\begin{enumerate}[label=\textbullet]
    \item Construire d'abord les symétriques de tous les sommets.
    \item Relier ensuite les points obtenus dans le même ordre.
\end{enumerate}

\begin{center}
\begin{tikzpicture}[scale=0.75]

% grille
\draw[step=1cm, gray!35, very thin] (-6,-4) grid (6,5);

% axe de symétrie horizontal
\draw[very thick, bleu] (-6,0)--(6,0);
\node[bleu, right] at (6,0) {$(d)$};

% figure initiale
\draw[very thick] (-5,2)--(-3,4)--(-1,3)--(-2,1)--(-4,1)--cycle;

% points
\filldraw (-5,2) circle (2pt) node[above left] {$A$};
\filldraw (-3,4) circle (2pt) node[above] {$B$};
\filldraw (-1,3) circle (2pt) node[above right] {$C$};
\filldraw (-2,1) circle (2pt) node[below right] {$D$};
\filldraw (-4,1) circle (2pt) node[below left] {$E$};

\end{tikzpicture}
\end{center}

\textbf{Questions :}

\begin{enumerate}[label=\arabic*)]
    \item Construire les points $A'$, $B'$, $C'$, $D'$ et $E'$. \points{3}
    \item Relier les points $A'$, $B'$, $C'$, $D'$ et $E'$ pour obtenir la figure symétrique. \points{2}
    \item La figure symétrique a-t-elle la même forme et les mêmes dimensions que la figure de départ ? Justifier. \points{1}
\end{enumerate}

\vspace{0.8cm}

Justification :

\dotfill

\dotfill

\dotfill

% ============================================================
\section*{Exercice 4 -- Construction avec le compas seul \points{5}}

Dans cet exercice, la droite $(d)$ est l'axe de symétrie de la figure complète.

\medskip

\textbf{Instrument imposé : compas seul.}

\medskip

On a placé deux points $M$ et $N$ sur l'axe de symétrie $(d)$.  
Tu dois construire les symétriques des points $A$, $B$ et $C$ en utilisant seulement le compas.

\medskip

\textbf{Indication :} si $A'$ est le symétrique de $A$ par rapport à $(d)$, alors
\[
MA = MA'
\qquad \text{et} \qquad
NA = NA'.
\]
Donc $A'$ est à l'intersection de deux arcs de cercle : l'un de centre $M$, l'autre de centre $N$.

\begin{center}
\begin{tikzpicture}[scale=0.9]

% axe oblique
\draw[very thick, bleu] (-5,-2)--(5,3);
\node[bleu, right] at (5,3) {$(d)$};

% points sur l'axe
\filldraw[bleu] (-2,-0.5) circle (2pt) node[below left] {$M$};
\filldraw[bleu] (2,1.5) circle (2pt) node[above right] {$N$};

% triangle
\draw[very thick] (-4,2)--(-1,3)--(-0.5,0.8)--cycle;
\filldraw (-4,2) circle (2pt) node[above left] {$A$};
\filldraw (-1,3) circle (2pt) node[above] {$B$};
\filldraw (-0.5,0.8) circle (2pt) node[right] {$C$};

\end{tikzpicture}
\end{center}

\begin{enumerate}[label=\arabic*)]
    \item Construire au compas seulement les points $A'$, $B'$ et $C'$. \points{3}
    \item Relier les points $A'$, $B'$ et $C'$ pour compléter la figure. \points{1}
    \item Pourquoi les points $M$ et $N$ ne changent-ils pas par symétrie par rapport à $(d)$ ? \points{1}
\end{enumerate}

\vspace{0.8cm}

Réponse :

\dotfill

\dotfill

\dotfill
\section*{Exercice 5 -- Compléter une figure par symétrie \points{5}}

La droite $(d)$ est un axe de symétrie de la figure complète.  
On a tracé seulement une moitié de la figure.

\medskip

\textbf{Compléter la figure pour qu'elle soit symétrique par rapport à la droite $(d)$.}

\begin{center}
\begin{tikzpicture}[scale=0.75]

% grille
\draw[step=1cm, gray!35, very thin] (-6,-4) grid (6,4);

% axe
\draw[very thick, bleu] (0,-4)--(0,4);
\node[bleu, above] at (0,4) {$(d)$};

% demi-figure côté gauche
\draw[very thick] (0,3)--(-2,2)--(-4,0)--(-3,-2)--(-1,-3)--(0,-2);

% points
\filldraw (0,3) circle (2pt) node[right] {$A$};
\filldraw (-2,2) circle (2pt) node[above left] {$B$};
\filldraw (-4,0) circle (2pt) node[left] {$C$};
\filldraw (-3,-2) circle (2pt) node[below left] {$D$};
\filldraw (-1,-3) circle (2pt) node[below] {$E$};
\filldraw (0,-2) circle (2pt) node[right] {$F$};

\end{tikzpicture}
\end{center}

\textbf{Questions :}

\begin{enumerate}[label=\arabic*)]
    \item Construire les symétriques des points $B$, $C$, $D$ et $E$. \points{3}
    \item Compléter la figure en reliant correctement les points. \points{1}
    \item Les points $A$ et $F$ ont-ils besoin d'être déplacés pour obtenir leurs symétriques ? Expliquer. \points{1}
\end{enumerate}

\vspace{0.8cm}

Réponse à la question 3 :

\dotfill

\dotfill

\dotfill


% ============================================================
\section*{Exercice 6 -- Construction avec règle et compas \points{6}}

On veut construire le symétrique du triangle $ABC$ par rapport à la droite $(d)$.

\medskip

\textbf{Instrument imposé : règle non graduée et compas.}

\medskip

\textbf{Méthode attendue :}
\begin{enumerate}[label=\textbullet]
    \item Pour construire le symétrique d'un point, on utilise des arcs de cercle centrés sur deux points de la droite $(d)$.
    \item On ne mesure pas les distances avec la règle graduée.
    \item On relie ensuite les points obtenus.
\end{enumerate}

\begin{center}
\begin{tikzpicture}[scale=0.75]

% cadre/grille légère
\draw[step=1cm, gray!20, very thin] (-6,-3) grid (6,5);

% axe oblique
\draw[very thick, bleu] (-5,-2)--(5,4);
\node[bleu, right] at (5,4) {$(d)$};

% deux points sur l'axe pour aider
\filldraw[bleu] (-2,-0.2) circle (2pt) node[below right] {$M$};
\filldraw[bleu] (2,2.2) circle (2pt) node[above left] {$N$};

% triangle
\draw[very thick] (-4,3)--(-2,4)--(-1,1)--cycle;
\filldraw (-4,3) circle (2pt) node[above left] {$A$};
\filldraw (-2,4) circle (2pt) node[above] {$B$};
\filldraw (-1,1) circle (2pt) node[below right] {$C$};

\end{tikzpicture}
\end{center}

\begin{enumerate}[label=\arabic*)]
    \item Construire les points $A'$, $B'$ et $C'$, symétriques des points $A$, $B$ et $C$ par rapport à $(d)$. \points{4}
    \item Tracer le triangle $A'B'C'$. \points{1}
    \item Expliquer pourquoi $ABC$ et $A'B'C'$ ont les mêmes dimensions (c'est à dire chaque cotés du triangle à les mêmes longueurs). \points{1}
\end{enumerate}

\vspace{0.6cm}

Explication :

\dotfill

\dotfill

\dotfill

\end{document}